\section{Supuestos del modelo}\label{sec:supuestos}

\subsection{El sistema en lazo cerrado}

Hasta aquí, las entradas $\mathbf{u}(t)$ se han tratado como forzamientos
exógenos. El sistema real, sin embargo, opera en \emph{lazo cerrado}: los
actuadores responden a las mediciones. Esta subsección formaliza esa
estructura antes de enunciar los supuestos del modelo.

\begin{definition}[Ley de control y consignas]\label{def41}
Sea $\mathbf{u}(t) = \bigl(Q(t), P_{\mathrm{cal}}(t), P_{\mathrm{hum}}(t)\bigr)$
el vector de entradas manipuladas ---caudal de ventilación, potencia de
calefacción y potencia de humidificación--- y
$\mathbf{r} = (T_{\mathrm{ref}}, H_{\mathrm{ref}}, \bar{c}_{\mathrm{CO_2}},
\bar{c}_{\mathrm{CH_4}})$ el vector de \emph{consignas} (valores de
referencia y límites de seguridad). Una \emph{ley de control} es un mapa
\begin{equation}
\mathbf{u}(t) = \pi\bigl(\mathbf{y}(t), \mathbf{x}_c(t); \mathbf{r}\bigr),
\label{eq:ley-control}
\end{equation}
donde $\mathbf{x}_c$ es el estado interno del controlador (por ejemplo, los
integradores de acciones PID). El \emph{sistema en lazo cerrado} es la
composición planta--controlador,
\begin{equation}
\dot{\mathbf{x}} = f\bigl(\mathbf{x}, \pi(g(\mathbf{x}), \mathbf{x}_c;
\mathbf{r}); \boldsymbol{\theta}\bigr),
\qquad
\dot{\mathbf{x}}_c = f_c\bigl(\mathbf{x}_c, g(\mathbf{x});
\mathbf{r}\bigr),
\label{eq:lazo-cerrado}
\end{equation}
con estado aumentado $(\mathbf{x}, \mathbf{x}_c)$ y exógenos reducidos a
las consignas $\mathbf{r}$ y las condiciones del \emph{inlet}.
\end{definition}

\begin{remark}[Estructura jerárquica del control]\label{rem:jerarquia}
El actuador $Q$ es compartido por todos los balances de gas: aumentar la
ventilación reduce CO$_2$ y CH$_4$ pero arrastra calor y humedad, en
conflicto con los lazos de $T$ y $H$. Se adopta por tanto una arquitectura
de dos niveles: lazos regulatorios de temperatura y humedad en el nivel
inferior, y una capa supervisora en el nivel superior que impone los
límites de seguridad de gases (incluidos H$_2$S y NH$_3$ como variables de
supervisión) y puede modificar las consignas o forzar ventilación máxima.
\end{remark}

\begin{remark}[Identificaci\'on en lazo cerrado]
\label{rem:id-lazo}
En~\eqref{eq:lazo-cerrado}, las entradas $u$ est\'an correlacionadas con las perturbaciones por construcci\'on: el controlador reacciona a ellas. Es la dificultad cl\'asica de la identificaci\'on en lazo cerrado~\citep{ljungSystemIdentification1999}: un estimador entrenado solo con datos de operaci\'on regulada puede aprender la ley $\pi$ en lugar de la fisiolog\'ia. En consecuencia, los datos de los cierres de ventilaci\'on (Observaci\'on~\ref{rem:cierres}) y de las ventanas con consignas en escal\'on desempe\~nar\'an el rol de datos de identificaci\'on, mientras que la operaci\'on regulada normal se reserva para validaci\'on (cf.~Secci\'on~6).
\end{remark}


\begin{definition}[Entrada de ventilaci\'on conmutada]
\label{def42}
Sea $\{t_k^c\}_{k \geq 1}$ el calendario de cierres programados (Observaci\'on~\ref{rem:cierres}) y $\Delta$ la duraci\'on de cada cierre. La ventilaci\'on opera en dos modos:
\begin{equation}
Q(t) =
\begin{cases}
Q_{\mathrm{op}}(t), & \text{modo ventilado: } t \notin \bigcup_k [t_k^c,\; t_k^c + \Delta],\\[4pt]
0, & \text{modo cerrado: } t \in [t_k^c,\; t_k^c + \Delta], \quad k = 1, 2, \ldots
\end{cases}
\label{eq:q-conmutada}
\end{equation}
donde en modo ventilado $Q_{\mathrm{op}}(t)$ est\'a determinado por la ley de control~\eqref{eq:ley-control}, y en modo cerrado la ventilaci\'on se anula por decisi\'on del protocolo, suspendi\'endose temporalmente la regulaci\'on.
\end{definition}

\begin{remark}[La conmutación no altera el estado]\label{rem:conmutacion}
A diferencia de un sistema híbrido con eventos de salto, en
\eqref{eq:q-conmutada} el estado $\mathbf{x}(t)$ es continuo en los
instantes de conmutación: lo que cambia es qué términos de los balances
están activos. En modo ventilado, los balances de gases y humedad incluyen
los flujos $J_{\mathrm{in}}$, $J_{\mathrm{out}}$ asociados a $Q$; en modo
cerrado esos términos se anulan y las concentraciones evolucionan solo por
producción interna,
\begin{equation}
\dot{c}_i = \frac{R_i(t)}{V_{\mathrm{aire}}},
\qquad t \in [t_k^{c},\, t_k^{c} + \Delta],
\label{eq:modo-cerrado}
\end{equation}
con $V_{\mathrm{aire}}$ el volumen de aire del contenedor. La
Ec.~\eqref{eq:modo-cerrado} es la base del método de cámara estática~\citep{hutchinsonImprovedSoilCover1981,ermolaevGreenhouseGasEmissions2019}: las
pendientes de acumulación miden directamente las tasas $R_i$
(Sección~\ref{sec:mediciones}).
\end{remark}

\begin{remark}[Duración del cierre]\label{rem:delta}
La duración $\Delta$ resuelve un compromiso entre detectabilidad y
seguridad. Por abajo, la acumulación debe superar con holgura la resolución
$\sigma_i$ de cada sensor:
\begin{equation}
\Delta \;\geq\; \max_i \frac{10\,\sigma_i\, V_{\mathrm{aire}}}{R_i^{\min}},
\label{eq:delta-min}
\end{equation}
evaluada en el período de menor actividad metabólica (inicio del ciclo).
Por arriba, el O$_2$ no debe caer por debajo de un piso seguro
$\underline{c}_{\mathrm{O_2}}$:
\begin{equation}
\Delta \;\leq\; \frac{\bigl(c_{\mathrm{O_2}}(t_k^{c})
- \underline{c}_{\mathrm{O_2}}\bigr)\, V_{\mathrm{aire}}}{R_{\mathrm{O_2}}^{\max}}.
\label{eq:delta-max}
\end{equation}
Ambas cotas son calculables \emph{antes} del experimento a partir del
modelo, y ajustables día a día por el gemelo digital conforme $R_i$
evoluciona (Sección~\ref{sec:estimacion}).
\end{remark}

\subsection{Supuestos de compartimentación}

\begin{description}
\item[A1. Mezcla homogénea.] Cada compartimento está uniformemente
distribuido en su volumen (Definición~\ref{def13}): un escalar describe su
estado y una medición puntual es representativa. Los gradientes locales en
la vecindad de la masa larval se aceptan como fuente de error de medición,
a vigilar en la validación.

\item[A2. Población homogénea.] Las $N(t)$ larvas comparten el estado
$(A, B, L)$ de una larva promedio (Definición~\ref{def31}); la mortalidad,
si ocurre, se admite mediante $N(t)$ decreciente.

\item[A3. Operación por lote.] Sin recargas de alimento ni de larvas en
$[0, t_f]$ (Observación~\ref{rem:batch}); el estado es continuo y las
únicas conmutaciones son las de ventilación (Definición~\ref{def42}).
\end{description}

\subsection{Supuestos fisiológicos}

\begin{description}
\item[B1. Asimilación limitante.] La asimilación de alimento es el proceso
metabólico limitante, $r_A = aB$ \citep{eriksenDynamicModellingFeed2022}, con $a$ modulada por
funciones logísticas del peso según el modelo de referencia.

\item[B2. Mantenimiento y costos de síntesis.] El mantenimiento es
proporcional a la biomasa estructural, $r_{\mathrm{CO_2},m} = mB$, y los
costos de síntesis $Y_B$, $Y_L$ son fracciones constantes de los flujos
respectivos (Ecs.~\eqref{eq:rco2m}--\eqref{eq:rco2l}).

\item[B3. Estado estacionario de asimilados.] El compartimento $A$
equilibra sus flujos en escalas rápidas frente a la dinámica de $B$ y $L$
(Ec.~1 de \citet{eriksenDynamicModellingFeed2022}), de modo que $\dot{A} \approx 0$ en
\eqref{eq:eriksen-edos}.

\item[B4. Parámetros constantes.] Se adoptan los parámetros fisiológicos
$(a, m, Y_B, Y_L, B_{\max}, \beta)$ como constantes durante el ciclo. Esta
es la práctica establecida en la literatura de referencia: el modelo de
\citet{eriksenDynamicModellingFeed2022} fue formulado y ajustado con parámetros constantes a
temperatura de cría controlada, con buen desempeño frente a datos de
distintos sustratos y humedades \citep{eriksenDynamicModellingFeed2022,eriksenMetabolicPerformanceFeed2024}. Dado que
el presente sistema mantiene $T$ regulada en la zona de confort larval
(Definición~\ref{def41}), la aproximación es razonable; la dependencia
explícita de los parámetros con $T$ (por ejemplo, factores de tipo
$Q_{10}$ o Arrhenius) queda planteada como extensión para el caso de
operación fuera de la zona regulada.
\end{description}